\section{Distinction loss, presentation, and continuation}\label{sec:barriers}

Fiber-based impossibility results depend only on an interface's distinction shadow. A restricted continuation family can additionally depend on the form in which the surviving distinctions are presented.

\subsection{Source resolution and aggregation are process structure}

A receiver interface may actually compute
\[
Q_j(H)=\bigl(L_j(H),E_j(H)\bigr),
\]
with transported exposure
\[
P_j:X\to\mathcal M_j,
\qquad
\operatorname{Agg}_j:\mathcal M_j\to Z_j^{\mathrm{exp}},
\qquad
E_j=\operatorname{Agg}_jP_j.
\]
For a finite source set, $P_j(H)=(m_{ji}(H))_{i\in S}$ is a common source-resolved state. These are distinct represented cuts, even though their distinction shadows can be compared. For example,
\[
P(x_1,x_2)=(x_1,x_2),
\qquad
\operatorname{Agg}(x_1,x_2)=x_1+x_2
\]
identifies $(1,-1)$ and $(0,0)$ after aggregation even though $x_1^2+x_2^2$ separates them.

\subsection{Fiber variation and deterministic barriers}

\begin{definition}[Target variation along a distinction shadow]\label{def:fiber-variation-v13}
For $T:X\to U$ and $f:X\to(W,d)$, define
\[
\Delta_T(f)
:=
\sup_{T(x)=T(x')}d(f(x),f(x')).
\]
\end{definition}

This depends only on $\mathsf D(T)$.

\begin{theorem}[Deterministic distinction barrier]\label{thm:deterministic-barrier-v13}
Let $E:X\twoheadrightarrow Z$ be a receiver interface and $G:Z\to W$ any continuation. Then
\[
\boxed{
\sup_{x\in X}d\bigl(f(x),G(E(x))\bigr)
\ge
\frac12\Delta_E(f).
}
\]
An exact continuation $f=GE$ exists if and only if $\Delta_E(f)=0$.
\end{theorem}

\begin{proof}
For $E(x)=E(x')=:z$, the triangle inequality gives
\[
d(f(x),f(x'))
\le d(f(x),G(z))+d(G(z),f(x')).
\]
Take the supremum over common-fiber pairs. Exact recovery is equivalent to fiber constancy by \Cref{lem:fiber-factorization-v13}.
\end{proof}

\begin{proposition}[Aggregation cannot improve target resolution]\label{prop:aggregation-resolution-v13}
If $E=\operatorname{Agg}P$, then
\[
\boxed{\Delta_E(f)\ge\Delta_P(f).}
\]
\end{proposition}

\begin{proof}
Every $P$-fiber is contained in an $E$-fiber.
\end{proof}

These exact statements are invariant under injective post-processing because it leaves the distinction shadow unchanged.

\subsection{Presented interfaces and continuation transport}

\begin{definition}[Attainable branch family]\label{def:attainable-branch-family-v13}
For $Q:X\twoheadrightarrow Z$ and
\[
\mathcal G\subseteq\{g:Z\to W\},
\]
define
\[
\boxed{
\mathfrak B(Q,\mathcal G)
:=\{gQ:g\in\mathcal G\}.
}
\]
\end{definition}

\begin{theorem}[Continuation transport]\label{thm:continuation-transport-v13}
Suppose $Q'=TQ$ for a bijection $T:Z\to Z'$, and let
\[
\mathcal G'\subseteq\{g':Z'\to W\},
\qquad
T^*\mathcal G':=\{g'T:g'\in\mathcal G'\}.
\]
Then
\[
\boxed{
\mathfrak B(Q',\mathcal G')
=
\mathfrak B(Q,T^*\mathcal G').
}
\]
Because $Q$ is surjective,
\[
\boxed{
\mathfrak B(Q,\mathcal G)
=
\mathfrak B(Q',\mathcal G')
\quad\Longleftrightarrow\quad
\mathcal G=T^*\mathcal G'.
}
\]
\end{theorem}

\begin{proof}
Every $g'Q'$ equals $g'TQ$. Moreover $g\mapsto gQ$ is injective on functions on $Z$ because $Q$ is surjective.
\end{proof}

Thus a bijective conversion preserves an attainable branch family only when the continuation class transports with it.

\begin{example}[Invertible affine mixing can change restricted accessibility]\label{ex:affine-presentation-v13}
Let $X=\R^2$,
\[
Q(u,v)=(u,v),
\qquad
Q'(u,v)=TQ(u,v)=(u+v,u-v).
\]
The linear map $T$ is invertible, so $Q\equiv_DQ'$. By \Cref{thm:affine-locality-v13}, however, both output coordinates of $T$ depend on both marked input coordinates.

Let the allowed scalar continuation class be exactly the one-coordinate functions
\[
\mathcal G_{\mathrm{coord}}
:=
\{\gamma\pi_1:\gamma:\R\to\R\}
\cup
\{\gamma\pi_2:\gamma:\R\to\R\}.
\]
For the target $f(u,v)=u$, we have $f=\pi_1Q$, so
\[
f\in\mathfrak B(Q,\mathcal G_{\mathrm{coord}}).
\]
But $f\notin\mathfrak B(Q',\mathcal G_{\mathrm{coord}})$. Indeed, the first coordinate of $Q'$ is equal at $(0,0)$ and $(1,-1)$ while $f$ differs there; the second coordinate is equal at $(0,0)$ and $(1,1)$ while $f$ differs there. Hence neither coordinate alone determines $u$. Exact recovery requires cross-coordinate mixing,
\[
\boxed{
u=\frac12(Q'_1+Q'_2).}
\]
\end{example}

The example isolates the point: invertibility preserves extensional distinctions, not the locality or admissibility of the computation needed to use a given presentation.

\subsection{Squared-loss distinction and continuation}

Let $E:X\to Z$ be measurable, let $H$ be $X$-valued, and write $Z_H:=E(H)$. Let $Y$ be finite-dimensional Euclidean-valued with finite second moment. For a nonempty family of measurable continuations $\mathcal G$ with square-integrable outputs, define
\[
m_E:=\E[Y\mid\sigma(Z_H)].
\]

\begin{theorem}[Squared-loss distinction/continuation decomposition]\label{thm:squared-decomposition-v13}
\[
\boxed{
\inf_{g\in\mathcal G}\E\|Y-g(Z_H)\|^2
=
\E\|Y-m_E\|^2
+
\inf_{g\in\mathcal G}\E\|m_E-g(Z_H)\|^2.
}
\]
\end{theorem}

\begin{proof}
For each $g$,
\[
Y-g(Z_H)=(Y-m_E)+(m_E-g(Z_H)),
\]
and the cross term has expectation zero by conditional expectation. Take the infimum.
\end{proof}

The first term is irreducible given the retained distinction structure; the second is the approximation defect of the selected continuation family on the actual presentation.

\begin{corollary}[Presentation invariance of the Bayes term]\label{cor:presentation-bayes-v13}
Suppose $E'=TE$ where $T$ is a measurable bijection on $\im E$ with measurable inverse. Then
\[
\sigma(E(H))=\sigma(E'(H)),
\]
so the Bayes term in \Cref{thm:squared-decomposition-v13} is identical for $E$ and $E'$. If continuation families satisfy
\[
\mathcal G=T^*\mathcal G',
\]
then their approximation terms and optimal squared losses are also equal. Without that transport condition, the approximation terms can differ despite identical distinction shadows.
\end{corollary}

\begin{proof}
The measurable inverse gives both sigma-algebra inclusions. The continuation claim follows from \Cref{thm:continuation-transport-v13}.
\end{proof}
